% Options for packages loaded elsewhere
\PassOptionsToPackage{unicode}{hyperref}
\PassOptionsToPackage{hyphens}{url}
\documentclass[
  11pt,
  a4paper,
]{article}
\usepackage{xcolor}
\usepackage[margin=18mm]{geometry}
\usepackage{amsmath,amssymb}
\setcounter{secnumdepth}{-\maxdimen} % remove section numbering
\usepackage{iftex}
\ifPDFTeX
  \usepackage[T1]{fontenc}
  \usepackage[utf8]{inputenc}
  \usepackage{textcomp} % provide euro and other symbols
\else % if luatex or xetex
  \usepackage{unicode-math} % this also loads fontspec
  \defaultfontfeatures{Scale=MatchLowercase}
  \defaultfontfeatures[\rmfamily]{Ligatures=TeX,Scale=1}
\fi
\ifPDFTeX\else
  % xetex/luatex font selection
    \setmainfont[]{Noto Serif}
    \setsansfont[]{Noto Sans}
    \setmonofont[]{Noto Sans Mono}
  \setmathfont[]{Noto Sans Math}
\fi
% Use upquote if available, for straight quotes in verbatim environments
\IfFileExists{upquote.sty}{\usepackage{upquote}}{}
\IfFileExists{microtype.sty}{% use microtype if available
  \usepackage[]{microtype}
  \UseMicrotypeSet[protrusion]{basicmath} % disable protrusion for tt fonts
}{}
\makeatletter
\@ifundefined{KOMAClassName}{% if non-KOMA class
  \IfFileExists{parskip.sty}{%
    \usepackage{parskip}
  }{% else
    \setlength{\parindent}{0pt}
    \setlength{\parskip}{6pt plus 2pt minus 1pt}}
}{% if KOMA class
  \KOMAoptions{parskip=half}}
\makeatother
\ifLuaTeX
\usepackage[bidi=basic]{babel}
\else
\usepackage[bidi=default]{babel}
\fi
\babelprovide[main,import]{russian}
\ifPDFTeX
\else
\babelfont{rm}[]{Noto Serif}
\fi
% get rid of language-specific shorthands (see #6817):
\let\LanguageShortHands\languageshorthands
\def\languageshorthands#1{}
\setlength{\emergencystretch}{3em} % prevent overfull lines
\providecommand{\tightlist}{%
  \setlength{\itemsep}{0pt}\setlength{\parskip}{0pt}}
\usepackage{bookmark}
\IfFileExists{xurl.sty}{\usepackage{xurl}}{} % add URL line breaks if available
\urlstyle{same}
\hypersetup{
  pdftitle={Билеты по дискретным случайным процессам},
  pdflang={ru-RU},
  hidelinks,
  pdfcreator={LaTeX via pandoc}}

\title{Билеты по дискретным случайным процессам}
\author{}
\date{}

\begin{document}
\maketitle

\section{Билет 08. Интегрируемость ограниченных измеримых функций.
Характеристические функции конечномерных
распределений}\label{ux431ux438ux43bux435ux442-08.-ux438ux43dux442ux435ux433ux440ux438ux440ux443ux435ux43cux43eux441ux442ux44c-ux43eux433ux440ux430ux43dux438ux447ux435ux43dux43dux44bux445-ux438ux437ux43cux435ux440ux438ux43cux44bux445-ux444ux443ux43dux43aux446ux438ux439.-ux445ux430ux440ux430ux43aux442ux435ux440ux438ux441ux442ux438ux447ux435ux441ux43aux438ux435-ux444ux443ux43dux43aux446ux438ux438-ux43aux43eux43dux435ux447ux43dux43eux43cux435ux440ux43dux44bux445-ux440ux430ux441ux43fux440ux435ux434ux435ux43bux435ux43dux438ux439}

Источники: Ширяев 1, гл. II, §§6,12; лекции ДСП.

\subsection{1. Короткий
ответ}\label{ux43aux43eux440ux43eux442ux43aux438ux439-ux43eux442ux432ux435ux442}

Если \((\Omega,\mathcal F,\mu)\) - пространство конечной меры, то всякая
ограниченная измеримая функция интегрируема:

\[
|f|\le M,\qquad \mu(\Omega)<\infty
\quad\Longrightarrow\quad
\int_\Omega |f|\,d\mu\le M\mu(\Omega)<\infty.
\]

В вероятностном пространстве это означает: всякая ограниченная случайная
величина имеет конечное математическое ожидание.

Характеристическая функция распределения - это его преобразование Фурье.
Для случайного вектора \(X=(X_1,\ldots,X_n)\):

\[
\varphi_X(t)
=\mathbb E e^{i\langle t,X\rangle}
=\int_{\mathbb R^n} e^{i\langle t,x\rangle}\,P_X(dx),
\qquad t\in\mathbb R^n.
\]

Она всегда существует, потому что

\[
\left|e^{i\langle t,X\rangle}\right|=1.
\]

Для случайного процесса конечномерные характеристические функции - это
характеристические функции векторов

\[
(X_{t_1},\ldots,X_{t_n}).
\]

Они однозначно задают соответствующие конечномерные распределения.

\subsection{2. Полный
ответ}\label{ux43fux43eux43bux43dux44bux439-ux43eux442ux432ux435ux442}

\subsubsection{2.1. Введение: основные идеи
билета}\label{ux432ux432ux435ux434ux435ux43dux438ux435-ux43eux441ux43dux43eux432ux43dux44bux435-ux438ux434ux435ux438-ux431ux438ux43bux435ux442ux430}

В билете соединяются две идеи. Первая относится к интегралу Лебега: на
пространстве конечной меры ограниченность и измеримость автоматически
дают интегрируемость. Вторая относится к распределениям:
характеристическая функция кодирует распределение как преобразование
Фурье и поэтому удобна для работы с суммами, предельными теоремами и
конечномерными распределениями случайных процессов.

\subsubsection{2.2. Интегрируемость ограниченных измеримых
функций}\label{ux438ux43dux442ux435ux433ux440ux438ux440ux443ux435ux43cux43eux441ux442ux44c-ux43eux433ux440ux430ux43dux438ux447ux435ux43dux43dux44bux445-ux438ux437ux43cux435ux440ux438ux43cux44bux445-ux444ux443ux43dux43aux446ux438ux439}

Пусть задано пространство с мерой

\[
(\Omega,\mathcal F,\mu).
\]

Функция \(f:\Omega\to\mathbb R\) или \(f:\Omega\to\mathbb C\) называется
интегрируемой, если

\[
\int_\Omega |f|\,d\mu<\infty.
\]

Для комплексной функции \(f=u+iv\) условие \(\int |f|\,d\mu<\infty\)
эквивалентно интегрируемости действительной и мнимой частей:
\(u,v\in L^1(\mu)\).

Если \(f\) измерима и ограничена, то существует \(M<\infty\) такое, что

\[
|f(\omega)|\le M
\]

для всех \(\omega\) или хотя бы почти всюду. Тогда

\[
\int_\Omega |f|\,d\mu
\le
\int_\Omega M\,d\mu
=M\mu(\Omega).
\]

Если \(\mu(\Omega)<\infty\), правая часть конечна, значит
\(f\in L^1(\mu)\). В вероятности \(\mu=P\) и \(P(\Omega)=1\), поэтому
оценка принимает особенно простой вид:

\[
\mathbb E|X|\le M.
\]

Главное ограничение этой теоремы: мера должна быть конечной. На
бесконечной мере утверждение неверно. Например, функция \(f(x)\equiv1\)
ограничена и измерима на \((\mathbb R,\mathcal B(\mathbb R),\lambda)\),
но

\[
\int_{\mathbb R} 1\,d\lambda=\infty.
\]

\subsubsection{2.3. Определение характеристической функции случайной
величины и
вектора}\label{ux43eux43fux440ux435ux434ux435ux43bux435ux43dux438ux435-ux445ux430ux440ux430ux43aux442ux435ux440ux438ux441ux442ux438ux447ux435ux441ux43aux43eux439-ux444ux443ux43dux43aux446ux438ux438-ux441ux43bux443ux447ux430ux439ux43dux43eux439-ux432ux435ux43bux438ux447ux438ux43dux44b-ux438-ux432ux435ux43aux442ux43eux440ux430}

Теперь перейдем к характеристическим функциям. Пусть \(X\) - случайная
величина с распределением \(P_X\) на \(\mathbb R\). Ее
характеристическая функция:

\[
\varphi_X(t)=\mathbb E e^{itX}
=\int_{\mathbb R} e^{itx}\,P_X(dx),
\qquad t\in\mathbb R.
\]

Если распределение задается функцией распределения \(F_X\), то пишут
также

\[
\varphi_X(t)=\int_{\mathbb R} e^{itx}\,dF_X(x).
\]

Если есть плотность \(p_X\), то

\[
\varphi_X(t)=\int_{\mathbb R} e^{itx}p_X(x)\,dx.
\]

Для случайного вектора \(X=(X_1,\ldots,X_n)\) характеристическая функция
задается формулой

\[
\varphi_X(t)
=\mathbb E e^{i\langle t,X\rangle}
=\int_{\mathbb R^n} e^{i\langle t,x\rangle}\,P_X(dx),
\qquad t=(t_1,\ldots,t_n)\in\mathbb R^n,
\]

где

\[
\langle t,x\rangle=t_1x_1+\cdots+t_nx_n.
\]

Эта величина всегда корректно определена: случайная величина
\(e^{i\langle t,X\rangle}\) комплекснозначна, измерима и ограничена по
модулю единицей. Поэтому она интегрируема на вероятностном пространстве.
Это важное преимущество характеристических функций перед производящими
функциями моментов: моментная производящая функция \(\mathbb E e^{tX}\)
может не существовать, а характеристическая функция существует всегда.

\subsubsection{2.4. Характеристические функции конечномерных
распределений
процесса}\label{ux445ux430ux440ux430ux43aux442ux435ux440ux438ux441ux442ux438ux447ux435ux441ux43aux438ux435-ux444ux443ux43dux43aux446ux438ux438-ux43aux43eux43dux435ux447ux43dux43eux43cux435ux440ux43dux44bux445-ux440ux430ux441ux43fux440ux435ux434ux435ux43bux435ux43dux438ux439-ux43fux440ux43eux446ux435ux441ux441ux430}

Для случайного процесса \((X_s)_{s\in T}\) конечномерное распределение
задается выбором конечного набора моментов времени

\[
s_1,\ldots,s_n\in T
\]

и законом случайного вектора

\[
(X_{s_1},\ldots,X_{s_n}).
\]

Его характеристическая функция:

\[
\varphi_{s_1,\ldots,s_n}(u_1,\ldots,u_n)
=
\mathbb E\exp\left(i\sum_{k=1}^n u_kX_{s_k}\right).
\]

Именно такие функции называются характеристическими функциями
конечномерных распределений процесса. Если знать их для всех конечных
наборов индексов, то знать все конечномерные распределения. В дальнейшем
это нужно для задания случайных процессов, гауссовских
последовательностей, стационарности и предельных теорем.

\subsubsection{2.5. Основные свойства и
непрерывность}\label{ux43eux441ux43dux43eux432ux43dux44bux435-ux441ux432ux43eux439ux441ux442ux432ux430-ux438-ux43dux435ux43fux440ux435ux440ux44bux432ux43dux43eux441ux442ux44c}

Основные свойства характеристической функции следуют из простых свойств
комплексной экспоненты и интеграла Лебега.

Во-первых,

\[
\varphi_X(0)=\mathbb E1=1.
\]

Во-вторых,

\[
|\varphi_X(t)|
=|\mathbb E e^{i\langle t,X\rangle}|
\le \mathbb E |e^{i\langle t,X\rangle}|
=1.
\]

В-третьих,

\[
\varphi_X(-t)=\overline{\varphi_X(t)}.
\]

В частности, для одномерной случайной величины с симметричным
относительно нуля распределением характеристическая функция
действительнозначна:

\[
\varphi_X(t)=\mathbb E\cos(tX).
\]

В-четвертых, характеристическая функция равномерно непрерывна. Для
одномерного случая:

\[
|\varphi(t+h)-\varphi(t)|
=
\left|\mathbb E\left(e^{itX}(e^{ihX}-1)\right)\right|
\le
\mathbb E|e^{ihX}-1|.
\]

При \(h\to0\) подынтегральное выражение стремится к нулю и ограничено
числом \(2\), поэтому по теореме Лебега о доминированной сходимости
правая часть стремится к нулю. Оценка не зависит от \(t\), значит
непрерывность равномерная.

\subsubsection{2.6. Положительная определенность и теорема
Бохнера}\label{ux43fux43eux43bux43eux436ux438ux442ux435ux43bux44cux43dux430ux44f-ux43eux43fux440ux435ux434ux435ux43bux435ux43dux43dux43eux441ux442ux44c-ux438-ux442ux435ux43eux440ux435ux43cux430-ux431ux43eux445ux43dux435ux440ux430}

В-пятых, характеристическая функция положительно определена:

\[
\sum_{j,k=1}^m c_j\overline{c_k}\,
\varphi_X(t_j-t_k)\ge0
\]

для любых \(t_1,\ldots,t_m\in\mathbb R^n\) и
\(c_1,\ldots,c_m\in\mathbb C\). Действительно,

\[
\sum_{j,k=1}^m c_j\overline{c_k}\,
\varphi_X(t_j-t_k)
=
\mathbb E\left|\sum_{j=1}^m c_je^{i\langle t_j,X\rangle}\right|^2
\ge0.
\]

Это свойство важно в обратную сторону: по теореме Бохнера непрерывная в
нуле положительно определенная функция \(\varphi\) с \(\varphi(0)=1\)
является характеристической функцией некоторого вероятностного
распределения. Для экзамена обычно достаточно знать необходимость этих
условий и смысл положительной определенности.

\subsubsection{2.7. Однозначность распределения и формула
обращения}\label{ux43eux434ux43dux43eux437ux43dux430ux447ux43dux43eux441ux442ux44c-ux440ux430ux441ux43fux440ux435ux434ux435ux43bux435ux43dux438ux44f-ux438-ux444ux43eux440ux43cux443ux43bux430-ux43eux431ux440ux430ux449ux435ux43dux438ux44f}

В-шестых, характеристическая функция однозначно определяет
распределение. Если две случайные величины или два случайных вектора
имеют одну и ту же характеристическую функцию, то их распределения
совпадают. В одномерном случае это можно усилить формулой обращения:
если \(a<b\) - точки непрерывности функции распределения \(F\), то

\[
F(b)-F(a)
=
\lim_{T\to\infty}
\frac1{2\pi}
\int_{-T}^{T}
\frac{e^{-ita}-e^{-itb}}{it}\,\varphi(t)\,dt.
\]

Если дополнительно \(\varphi\in L^1(\mathbb R)\), то распределение имеет
непрерывную плотность

\[
f(x)=\frac1{2\pi}\int_{\mathbb R} e^{-itx}\varphi(t)\,dt.
\]

Смысл: характеристическая функция - не просто вспомогательная величина,
а полная запись закона распределения.

\subsubsection{2.8. Характеристические функции независимых величин и
преобразований}\label{ux445ux430ux440ux430ux43aux442ux435ux440ux438ux441ux442ux438ux447ux435ux441ux43aux438ux435-ux444ux443ux43dux43aux446ux438ux438-ux43dux435ux437ux430ux432ux438ux441ux438ux43cux44bux445-ux432ux435ux43bux438ux447ux438ux43d-ux438-ux43fux440ux435ux43eux431ux440ux430ux437ux43eux432ux430ux43dux438ux439}

В-седьмых, характеристические функции особенно удобны для сумм
независимых величин. Если \(X\) и \(Y\) независимы, то

\[
\varphi_{X+Y}(t)
=\mathbb E e^{it(X+Y)}
=\mathbb E(e^{itX}e^{itY})
=\mathbb E e^{itX}\,\mathbb E e^{itY}
=\varphi_X(t)\varphi_Y(t).
\]

Для независимых \(X_1,\ldots,X_n\):

\[
\varphi_{X_1+\cdots+X_n}(t)
=
\prod_{k=1}^n\varphi_{X_k}(t).
\]

Именно поэтому характеристические функции используются в центральной
предельной теореме: свертка распределений превращается в произведение
функций.

Для векторов действует аналогичная формула. Если \(X\) и \(Y\)
независимые случайные векторы в \(\mathbb R^n\), то

\[
\varphi_{X+Y}(t)=\varphi_X(t)\varphi_Y(t).
\]

А если компоненты \(X_1,\ldots,X_n\) независимы, то

\[
\varphi_{(X_1,\ldots,X_n)}(t_1,\ldots,t_n)
=
\prod_{k=1}^n\varphi_{X_k}(t_k).
\]

Обратное тоже верно в стандартной форме: если совместная
характеристическая функция вектора факторизуется в произведение
одномерных характеристических функций координат, то координаты
независимы, потому что характеристическая функция однозначно задает
совместное распределение.

Еще одно полезное свойство - поведение при линейных преобразованиях.
Если \(Y=AX+b\), где \(A\) - матрица, \(b\) - вектор, то

\[
\varphi_Y(t)
=e^{i\langle t,b\rangle}\varphi_X(A^Tt).
\]

В одномерном случае для \(Y=aX+b\):

\[
\varphi_Y(t)=e^{itb}\varphi_X(at).
\]

Для конечномерных распределений процесса это означает, что линейные
комбинации координат также удобно описываются через одну и ту же
характеристическую функцию:

\[
\mathbb E e^{i(u_1X_{s_1}+\cdots+u_nX_{s_n})}.
\]

\subsubsection{2.9.
Заключение}\label{ux437ux430ux43aux43bux44eux447ux435ux43dux438ux435}

Итак, в этом билете главная логика такая: ограниченные измеримые функции
на конечной мере можно интегрировать; комплексная экспонента от
случайного вектора всегда ограничена; поэтому характеристическая функция
всегда существует и является надежным способом кодировать конечномерные
распределения.

\subsection{3. Основные
определения}\label{ux43eux441ux43dux43eux432ux43dux44bux435-ux43eux43fux440ux435ux434ux435ux43bux435ux43dux438ux44f}

\subsubsection{Интегрируемая
функция}\label{ux438ux43dux442ux435ux433ux440ux438ux440ux443ux435ux43cux430ux44f-ux444ux443ux43dux43aux446ux438ux44f}

Измеримая функция \(f\) называется интегрируемой относительно меры
\(\mu\), если

\[
\int |f|\,d\mu<\infty.
\]

В этом случае пишут \(f\in L^1(\mu)\).

\subsubsection{Ограниченная
функция}\label{ux43eux433ux440ux430ux43dux438ux447ux435ux43dux43dux430ux44f-ux444ux443ux43dux43aux446ux438ux44f}

Функция \(f\) ограничена, если существует \(M<\infty\) такое, что

\[
|f(\omega)|\le M
\]

для всех \(\omega\) или почти всюду. Для интеграла Лебега достаточно
ограничения почти всюду.

\subsubsection{Конечная
мера}\label{ux43aux43eux43dux435ux447ux43dux430ux44f-ux43cux435ux440ux430}

Мера \(\mu\) называется конечной, если

\[
\mu(\Omega)<\infty.
\]

Вероятностная мера всегда конечна, потому что \(P(\Omega)=1\).

\subsubsection{Комплекснозначная случайная
величина}\label{ux43aux43eux43cux43fux43bux435ux43aux441ux43dux43eux437ux43dux430ux447ux43dux430ux44f-ux441ux43bux443ux447ux430ux439ux43dux430ux44f-ux432ux435ux43bux438ux447ux438ux43dux430}

Случайная величина \(Z=U+iV\), где \(U\) и \(V\) - действительные
случайные величины. Ее математическое ожидание задается как

\[
\mathbb EZ=\mathbb EU+i\mathbb EV,
\]

если соответствующие ожидания существуют. Если \(\mathbb E|Z|<\infty\),
то ожидание существует.

\subsubsection{Распределение случайного
вектора}\label{ux440ux430ux441ux43fux440ux435ux434ux435ux43bux435ux43dux438ux435-ux441ux43bux443ux447ux430ux439ux43dux43eux433ux43e-ux432ux435ux43aux442ux43eux440ux430}

Для \(X:\Omega\to\mathbb R^n\) распределение \(P_X\) задается формулой

\[
P_X(B)=P\{X\in B\},
\qquad B\in\mathcal B(\mathbb R^n).
\]

\subsubsection{\texorpdfstring{Характеристическая функция распределения
на
\(\mathbb R^n\)}{Характеристическая функция распределения на \textbackslash mathbb R\^{}n}}\label{ux445ux430ux440ux430ux43aux442ux435ux440ux438ux441ux442ux438ux447ux435ux441ux43aux430ux44f-ux444ux443ux43dux43aux446ux438ux44f-ux440ux430ux441ux43fux440ux435ux434ux435ux43bux435ux43dux438ux44f-ux43dux430-mathbb-rn}

Если \(\mu\) - вероятностная мера на \(\mathbb R^n\), то

\[
\varphi_\mu(t)=\int_{\mathbb R^n} e^{i\langle t,x\rangle}\,\mu(dx),
\qquad t\in\mathbb R^n.
\]

\subsubsection{Характеристическая функция случайного
вектора}\label{ux445ux430ux440ux430ux43aux442ux435ux440ux438ux441ux442ux438ux447ux435ux441ux43aux430ux44f-ux444ux443ux43dux43aux446ux438ux44f-ux441ux43bux443ux447ux430ux439ux43dux43eux433ux43e-ux432ux435ux43aux442ux43eux440ux430}

Для \(X=(X_1,\ldots,X_n)\):

\[
\varphi_X(t)=\mathbb E e^{i\langle t,X\rangle}.
\]

\subsubsection{Характеристическая функция случайной
величины}\label{ux445ux430ux440ux430ux43aux442ux435ux440ux438ux441ux442ux438ux447ux435ux441ux43aux430ux44f-ux444ux443ux43dux43aux446ux438ux44f-ux441ux43bux443ux447ux430ux439ux43dux43eux439-ux432ux435ux43bux438ux447ux438ux43dux44b}

Частный случай при \(n=1\):

\[
\varphi_X(t)=\mathbb E e^{itX}.
\]

\subsubsection{Конечномерное распределение
процесса}\label{ux43aux43eux43dux435ux447ux43dux43eux43cux435ux440ux43dux43eux435-ux440ux430ux441ux43fux440ux435ux434ux435ux43bux435ux43dux438ux435-ux43fux440ux43eux446ux435ux441ux441ux430}

Для процесса \((X_s)_{s\in T}\) и индексов \(s_1,\ldots,s_n\) это
распределение вектора

\[
(X_{s_1},\ldots,X_{s_n}).
\]

\subsubsection{Характеристическая функция конечномерного
распределения}\label{ux445ux430ux440ux430ux43aux442ux435ux440ux438ux441ux442ux438ux447ux435ux441ux43aux430ux44f-ux444ux443ux43dux43aux446ux438ux44f-ux43aux43eux43dux435ux447ux43dux43eux43cux435ux440ux43dux43eux433ux43e-ux440ux430ux441ux43fux440ux435ux434ux435ux43bux435ux43dux438ux44f}

\[
\varphi_{s_1,\ldots,s_n}(u_1,\ldots,u_n)
=
\mathbb E\exp\left(i\sum_{k=1}^n u_kX_{s_k}\right).
\]

\subsubsection{Положительно определенная
функция}\label{ux43fux43eux43bux43eux436ux438ux442ux435ux43bux44cux43dux43e-ux43eux43fux440ux435ux434ux435ux43bux435ux43dux43dux430ux44f-ux444ux443ux43dux43aux446ux438ux44f}

Функция \(\varphi:\mathbb R^n\to\mathbb C\) положительно определена,
если для любых \(t_1,\ldots,t_m\in\mathbb R^n\) и
\(c_1,\ldots,c_m\in\mathbb C\)

\[
\sum_{j,k=1}^m c_j\overline{c_k}\varphi(t_j-t_k)\ge0.
\]

\subsection{4. Основные теоремы и
свойства}\label{ux43eux441ux43dux43eux432ux43dux44bux435-ux442ux435ux43eux440ux435ux43cux44b-ux438-ux441ux432ux43eux439ux441ux442ux432ux430}

\subsubsection{Теорема 1. Ограниченная измеримая функция интегрируема на
конечной
мере}\label{ux442ux435ux43eux440ux435ux43cux430-1.-ux43eux433ux440ux430ux43dux438ux447ux435ux43dux43dux430ux44f-ux438ux437ux43cux435ux440ux438ux43cux430ux44f-ux444ux443ux43dux43aux446ux438ux44f-ux438ux43dux442ux435ux433ux440ux438ux440ux443ux435ux43cux430-ux43dux430-ux43aux43eux43dux435ux447ux43dux43eux439-ux43cux435ux440ux435}

Пусть \((\Omega,\mathcal F,\mu)\) - пространство конечной меры, \(f\)
измерима и \(|f|\le M\) почти всюду. Тогда \(f\in L^1(\mu)\) и

\[
\int |f|\,d\mu\le M\mu(\Omega).
\]

\subsubsection{Доказательство}\label{ux434ux43eux43aux430ux437ux430ux442ux435ux43bux44cux441ux442ux432ux43e}

Так как \(|f|\le M\) почти всюду, по монотонности интеграла

\[
\int |f|\,d\mu\le\int M\,d\mu=M\mu(\Omega)<\infty.
\]

Значит \(f\) интегрируема.

\subsubsection{Следствие для
вероятности}\label{ux441ux43bux435ux434ux441ux442ux432ux438ux435-ux434ux43bux44f-ux432ux435ux440ux43eux44fux442ux43dux43eux441ux442ux438}

Если \(X\) - ограниченная случайная величина, то

\[
\mathbb E|X|\le M.
\]

Поэтому \(X\) имеет математическое ожидание.

\subsubsection{Следствие для характеристической
функции}\label{ux441ux43bux435ux434ux441ux442ux432ux438ux435-ux434ux43bux44f-ux445ux430ux440ux430ux43aux442ux435ux440ux438ux441ux442ux438ux447ux435ux441ux43aux43eux439-ux444ux443ux43dux43aux446ux438ux438}

Для любого случайного вектора \(X\) и любого \(t\in\mathbb R^n\)
случайная величина

\[
e^{i\langle t,X\rangle}
\]

ограничена по модулю единицей, значит интегрируема. Поэтому
\(\varphi_X(t)\) существует для всех \(t\).

\subsubsection{Свойство 1.
Нормировка}\label{ux441ux432ux43eux439ux441ux442ux432ux43e-1.-ux43dux43eux440ux43cux438ux440ux43eux432ux43aux430}

\[
\varphi_X(0)=1.
\]

\subsubsection{Свойство 2.
Ограниченность}\label{ux441ux432ux43eux439ux441ux442ux432ux43e-2.-ux43eux433ux440ux430ux43dux438ux447ux435ux43dux43dux43eux441ux442ux44c}

\[
|\varphi_X(t)|\le1.
\]

\subsubsection{Свойство 3. Сопряженная
симметрия}\label{ux441ux432ux43eux439ux441ux442ux432ux43e-3.-ux441ux43eux43fux440ux44fux436ux435ux43dux43dux430ux44f-ux441ux438ux43cux43cux435ux442ux440ux438ux44f}

\[
\varphi_X(-t)=\overline{\varphi_X(t)}.
\]

\subsubsection{Свойство 4. Равномерная
непрерывность}\label{ux441ux432ux43eux439ux441ux442ux432ux43e-4.-ux440ux430ux432ux43dux43eux43cux435ux440ux43dux430ux44f-ux43dux435ux43fux440ux435ux440ux44bux432ux43dux43eux441ux442ux44c}

Характеристическая функция равномерно непрерывна на \(\mathbb R^n\).

\subsubsection{Идея
доказательства}\label{ux438ux434ux435ux44f-ux434ux43eux43aux430ux437ux430ux442ux435ux43bux44cux441ux442ux432ux430}

Для \(h\to0\):

\[
|\varphi(t+h)-\varphi(t)|
\le
\mathbb E|e^{i\langle h,X\rangle}-1|
\to0
\]

по мажорантной теореме Лебега, так как модуль подынтегрального выражения
не больше \(2\).

\subsubsection{Свойство 5. Положительная
определенность}\label{ux441ux432ux43eux439ux441ux442ux432ux43e-5.-ux43fux43eux43bux43eux436ux438ux442ux435ux43bux44cux43dux430ux44f-ux43eux43fux440ux435ux434ux435ux43bux435ux43dux43dux43eux441ux442ux44c}

Характеристическая функция положительно определена:

\[
\sum_{j,k=1}^m c_j\overline{c_k}\varphi(t_j-t_k)\ge0.
\]

\subsubsection{Доказательство}\label{ux434ux43eux43aux430ux437ux430ux442ux435ux43bux44cux441ux442ux432ux43e-1}

\[
\sum_{j,k=1}^m c_j\overline{c_k}\varphi(t_j-t_k)
=
\mathbb E\left|\sum_{j=1}^m c_je^{i\langle t_j,X\rangle}\right|^2
\ge0.
\]

\subsubsection{Свойство 6.
Единственность}\label{ux441ux432ux43eux439ux441ux442ux432ux43e-6.-ux435ux434ux438ux43dux441ux442ux432ux435ux43dux43dux43eux441ux442ux44c}

Если

\[
\varphi_X(t)=\varphi_Y(t)
\qquad \text{для всех }t,
\]

то

\[
P_X=P_Y.
\]

То есть характеристическая функция однозначно определяет распределение.

\subsubsection{Свойство 7. Формула обращения в одномерном
случае}\label{ux441ux432ux43eux439ux441ux442ux432ux43e-7.-ux444ux43eux440ux43cux443ux43bux430-ux43eux431ux440ux430ux449ux435ux43dux438ux44f-ux432-ux43eux434ux43dux43eux43cux435ux440ux43dux43eux43c-ux441ux43bux443ux447ux430ux435}

Если \(a<b\) - точки непрерывности функции распределения \(F\), то

\[
F(b)-F(a)
=
\lim_{T\to\infty}
\frac1{2\pi}
\int_{-T}^{T}
\frac{e^{-ita}-e^{-itb}}{it}\varphi(t)\,dt.
\]

Если \(\varphi\in L^1(\mathbb R)\), то есть плотность

\[
f(x)=\frac1{2\pi}\int_{\mathbb R} e^{-itx}\varphi(t)\,dt.
\]

\subsubsection{Свойство 8. Линейное
преобразование}\label{ux441ux432ux43eux439ux441ux442ux432ux43e-8.-ux43bux438ux43dux435ux439ux43dux43eux435-ux43fux440ux435ux43eux431ux440ux430ux437ux43eux432ux430ux43dux438ux435}

Если \(Y=AX+b\), то

\[
\varphi_Y(t)=e^{i\langle t,b\rangle}\varphi_X(A^Tt).
\]

В частности, если \(Y=aX+b\), то

\[
\varphi_Y(t)=e^{itb}\varphi_X(at).
\]

\subsubsection{Свойство 9. Сумма независимых
величин}\label{ux441ux432ux43eux439ux441ux442ux432ux43e-9.-ux441ux443ux43cux43cux430-ux43dux435ux437ux430ux432ux438ux441ux438ux43cux44bux445-ux432ux435ux43bux438ux447ux438ux43d}

Если \(X\) и \(Y\) независимы, то

\[
\varphi_{X+Y}(t)=\varphi_X(t)\varphi_Y(t).
\]

Для независимых \(X_1,\ldots,X_n\):

\[
\varphi_{X_1+\cdots+X_n}(t)=\prod_{k=1}^n\varphi_{X_k}(t).
\]

\subsubsection{Свойство 10. Независимость координат через
факторизацию}\label{ux441ux432ux43eux439ux441ux442ux432ux43e-10.-ux43dux435ux437ux430ux432ux438ux441ux438ux43cux43eux441ux442ux44c-ux43aux43eux43eux440ux434ux438ux43dux430ux442-ux447ux435ux440ux435ux437-ux444ux430ux43aux442ux43eux440ux438ux437ux430ux446ux438ux44e}

Координаты \(X_1,\ldots,X_n\) независимы тогда и только тогда, когда

\[
\varphi_{(X_1,\ldots,X_n)}(t_1,\ldots,t_n)
=
\prod_{k=1}^n\varphi_{X_k}(t_k).
\]

\subsubsection{Свойство 11. Моменты через
производные}\label{ux441ux432ux43eux439ux441ux442ux432ux43e-11.-ux43cux43eux43cux435ux43dux442ux44b-ux447ux435ux440ux435ux437-ux43fux440ux43eux438ux437ux432ux43eux434ux43dux44bux435}

Если \(\mathbb E|X|^m<\infty\), то характеристическая функция имеет
производные до порядка \(m\), и

\[
\varphi_X^{(k)}(0)=i^k\mathbb E X^k,
\qquad k\le m.
\]

Эта связь подробно разбирается в Билете 13. В этом билете важно не
перепутать: характеристическая функция существует всегда, а моменты и
производные нужного порядка могут не существовать.

\subsubsection{Свойство 12. Согласованность конечномерных
характеристических
функций}\label{ux441ux432ux43eux439ux441ux442ux432ux43e-12.-ux441ux43eux433ux43bux430ux441ux43eux432ux430ux43dux43dux43eux441ux442ux44c-ux43aux43eux43dux435ux447ux43dux43eux43cux435ux440ux43dux44bux445-ux445ux430ux440ux430ux43aux442ux435ux440ux438ux441ux442ux438ux447ux435ux441ux43aux438ux445-ux444ux443ux43dux43aux446ux438ux439}

Если конечномерные распределения процесса согласованы, то их
характеристические функции тоже согласованы. Например,

\[
\varphi_{s_1,\ldots,s_n,s_{n+1}}(u_1,\ldots,u_n,0)
=
\varphi_{s_1,\ldots,s_n}(u_1,\ldots,u_n).
\]

Перестановка индексов \(s_k\) соответствует такой же перестановке
аргументов \(u_k\). Это характеристическая версия согласованности из
Билета 15.

\subsection{5. Примеры}\label{ux43fux440ux438ux43cux435ux440ux44b}

\subsubsection{Пример 1. Ограниченная случайная
величина}\label{ux43fux440ux438ux43cux435ux440-1.-ux43eux433ux440ux430ux43dux438ux447ux435ux43dux43dux430ux44f-ux441ux43bux443ux447ux430ux439ux43dux430ux44f-ux432ux435ux43bux438ux447ux438ux43dux430}

Пусть \(X\) принимает значения в отрезке \([-3,3]\). Тогда

\[
|X|\le3
\]

и

\[
\mathbb E|X|\le3.
\]

Значит \(X\) интегрируема, даже если мы не знаем ее точного
распределения.

\subsubsection{Пример 2. Почему нужна конечность
меры}\label{ux43fux440ux438ux43cux435ux440-2.-ux43fux43eux447ux435ux43cux443-ux43dux443ux436ux43dux430-ux43aux43eux43dux435ux447ux43dux43eux441ux442ux44c-ux43cux435ux440ux44b}

На \((\mathbb R,\mathcal B(\mathbb R),\lambda)\) функция

\[
f(x)\equiv1
\]

ограничена и измерима, но

\[
\int_{\mathbb R} |f(x)|\,dx=\infty.
\]

Значит ограниченность сама по себе не гарантирует интегрируемость на
бесконечной мере.

\subsubsection{Пример 3. Вырожденное
распределение}\label{ux43fux440ux438ux43cux435ux440-3.-ux432ux44bux440ux43eux436ux434ux435ux43dux43dux43eux435-ux440ux430ux441ux43fux440ux435ux434ux435ux43bux435ux43dux438ux435}

Если \(X=c\) почти наверное, то

\[
\varphi_X(t)=\mathbb E e^{itc}=e^{itc}.
\]

Это характеристическая функция меры Дирака \(\delta_c\).

\subsubsection{Пример 4. Бернуллиевская случайная
величина}\label{ux43fux440ux438ux43cux435ux440-4.-ux431ux435ux440ux43dux443ux43bux43bux438ux435ux432ux441ux43aux430ux44f-ux441ux43bux443ux447ux430ux439ux43dux430ux44f-ux432ux435ux43bux438ux447ux438ux43dux430}

Пусть

\[
P\{X=1\}=p,\qquad P\{X=0\}=q,\qquad p+q=1.
\]

Тогда

\[
\varphi_X(t)=\mathbb E e^{itX}=q+pe^{it}.
\]

Если \(X_1,\ldots,X_n\) независимы и имеют такое же распределение, то
для \(S_n=X_1+\cdots+X_n\)

\[
\varphi_{S_n}(t)=(q+pe^{it})^n.
\]

Это характеристическая функция биномиального распределения.

\subsubsection{Пример 5. Дискретный
вектор}\label{ux43fux440ux438ux43cux435ux440-5.-ux434ux438ux441ux43aux440ux435ux442ux43dux44bux439-ux432ux435ux43aux442ux43eux440}

Пусть случайный вектор \(X\) принимает значения
\(x^{(1)},\ldots,x^{(m)}\in\mathbb R^n\) с вероятностями
\(p_1,\ldots,p_m\). Тогда

\[
\varphi_X(t)=\sum_{r=1}^m p_r e^{i\langle t,x^{(r)}\rangle}.
\]

Например, если \(X=(X_1,X_2)\) принимает значения \((0,0),(1,0),(0,1)\)
с вероятностями \(p_0,p_1,p_2\), то

\[
\varphi_X(t_1,t_2)=p_0+p_1e^{it_1}+p_2e^{it_2}.
\]

\subsubsection{Пример 6. Нормальное
распределение}\label{ux43fux440ux438ux43cux435ux440-6.-ux43dux43eux440ux43cux430ux43bux44cux43dux43eux435-ux440ux430ux441ux43fux440ux435ux434ux435ux43bux435ux43dux438ux435}

Если

\[
X\sim N(m,\sigma^2),
\]

то

\[
\varphi_X(t)=\exp\left(imt-\frac{\sigma^2t^2}{2}\right).
\]

Для гауссовского вектора \(X\sim N(m,\Sigma)\):

\[
\varphi_X(t)
=
\exp\left(i\langle t,m\rangle-\frac12 t^T\Sigma t\right).
\]

Эта формула будет основной в Билете 18.

\subsubsection{Пример 7. Независимые
координаты}\label{ux43fux440ux438ux43cux435ux440-7.-ux43dux435ux437ux430ux432ux438ux441ux438ux43cux44bux435-ux43aux43eux43eux440ux434ux438ux43dux430ux442ux44b}

Пусть \(X_1\) и \(X_2\) независимы. Тогда

\[
\varphi_{(X_1,X_2)}(t_1,t_2)
=
\mathbb E e^{i(t_1X_1+t_2X_2)}
=
\varphi_{X_1}(t_1)\varphi_{X_2}(t_2).
\]

Если такая факторизация нарушается, координаты не являются независимыми.

\subsubsection{Пример 8. Характеристическая функция существует без
математического
ожидания}\label{ux43fux440ux438ux43cux435ux440-8.-ux445ux430ux440ux430ux43aux442ux435ux440ux438ux441ux442ux438ux447ux435ux441ux43aux430ux44f-ux444ux443ux43dux43aux446ux438ux44f-ux441ux443ux449ux435ux441ux442ux432ux443ux435ux442-ux431ux435ux437-ux43cux430ux442ux435ux43cux430ux442ux438ux447ux435ux441ux43aux43eux433ux43e-ux43eux436ux438ux434ux430ux43dux438ux44f}

У распределения Коши с плотностью

\[
p(x)=\frac1{\pi(1+x^2)}
\]

не существует \(\mathbb E|X|\), но характеристическая функция существует
и равна

\[
\varphi_X(t)=e^{-|t|}.
\]

Это хороший пример против ошибки ``если нет матожидания, то нет
характеристической функции''.

\subsubsection{Пример 9. Конечномерная характеристическая функция
процесса}\label{ux43fux440ux438ux43cux435ux440-9.-ux43aux43eux43dux435ux447ux43dux43eux43cux435ux440ux43dux430ux44f-ux445ux430ux440ux430ux43aux442ux435ux440ux438ux441ux442ux438ux447ux435ux441ux43aux430ux44f-ux444ux443ux43dux43aux446ux438ux44f-ux43fux440ux43eux446ux435ux441ux441ux430}

Пусть \((X_k)_{k\ge1}\) - случайная последовательность. Для трех
моментов времени \(1,2,5\):

\[
\varphi_{1,2,5}(u_1,u_2,u_3)
=
\mathbb E e^{i(u_1X_1+u_2X_2+u_3X_5)}.
\]

Если \(X_1,X_2,X_5\) независимы, то

\[
\varphi_{1,2,5}(u_1,u_2,u_3)
=
\varphi_{X_1}(u_1)\varphi_{X_2}(u_2)\varphi_{X_5}(u_3).
\]

Если последовательность стационарна в узком смысле, то такие функции не
меняются при общем сдвиге индексов:

\[
\varphi_{1,2,5}=\varphi_{1+h,2+h,5+h}.
\]

\subsection{6. Схемы доказательств /
обоснований}\label{ux441ux445ux435ux43cux44b-ux434ux43eux43aux430ux437ux430ux442ux435ux43bux44cux441ux442ux432-ux43eux431ux43eux441ux43dux43eux432ux430ux43dux438ux439}

\textbf{Почему ограниченная измеримая функция интегрируема на конечной
мере.}

Достаточно сравнить \(|f|\) с константой \(M\):

\[
0\le |f|\le M.
\]

По монотонности интеграла

\[
\int |f|\,d\mu\le \int M\,d\mu=M\mu(\Omega).
\]

Конечность \(\mu(\Omega)\) завершает доказательство.

\textbf{Почему характеристическая функция всегда существует.}

Для любого \(t\):

\[
\left|e^{i\langle t,X\rangle}\right|=1.
\]

Это ограниченная измеримая комплекснозначная случайная величина. На
вероятностном пространстве любая ограниченная случайная величина
интегрируема, значит

\[
\mathbb E e^{i\langle t,X\rangle}
\]

корректно определено.

\textbf{Почему \(|\varphi(t)|\le1\).}

Используется неравенство

\[
|\mathbb EZ|\le\mathbb E|Z|
\]

для интегрируемой комплексной случайной величины \(Z\). При
\(Z=e^{i\langle t,X\rangle}\) получаем

\[
|\varphi(t)|\le \mathbb E1=1.
\]

\textbf{Почему характеристическая функция непрерывна.}

Разность:

\[
\varphi(t+h)-\varphi(t)
=
\mathbb E\left[e^{i\langle t,X\rangle}(e^{i\langle h,X\rangle}-1)\right].
\]

Подынтегральное выражение по модулю не больше \(2\) и стремится к нулю
при \(h\to0\). Поэтому можно применить теорему Лебега о доминированной
сходимости.

\textbf{Почему характеристическая функция суммы независимых величин
равна произведению.}

Если \(X\) и \(Y\) независимы, то \(e^{itX}\) и \(e^{itY}\) тоже
независимы как измеримые функции от независимых случайных величин. Тогда

\[
\mathbb E(e^{itX}e^{itY})
=
\mathbb E e^{itX}\,\mathbb E e^{itY}.
\]

\textbf{Почему характеристическая функция задает распределение.}

Это теорема единственности. Идея доказательства: по \(\varphi(t)\) можно
восстановить интегралы от достаточно хороших функций через обратное
преобразование Фурье. Затем индикаторы интервалов аппроксимируются
такими функциями, и восстанавливаются значения

\[
P\{a<X\le b\}=F(b)-F(a)
\]

в точках непрерывности \(F\). Следовательно, восстанавливается вся
функция распределения.

\textbf{Почему конечномерные характеристические функции задают
конечномерные распределения процесса.}

Для каждого набора \(s_1,\ldots,s_n\) функция

\[
\varphi_{s_1,\ldots,s_n}
\]

является характеристической функцией случайного вектора
\((X_{s_1},\ldots,X_{s_n})\). По теореме единственности она задает его
распределение. Если это известно для всех конечных наборов индексов, то
известны все конечномерные распределения процесса.

\subsection{7. Что написать на
доске}\label{ux447ux442ux43e-ux43dux430ux43fux438ux441ux430ux442ux44c-ux43dux430-ux434ux43eux441ux43aux435}

\begin{enumerate}
\def\labelenumi{\arabic{enumi}.}
\tightlist
\item
  Интегрируемость ограниченной функции:
\end{enumerate}

\[
|f|\le M,\quad \mu(\Omega)<\infty
\quad\Rightarrow\quad
\int |f|\,d\mu\le M\mu(\Omega)<\infty.
\]

\begin{enumerate}
\def\labelenumi{\arabic{enumi}.}
\setcounter{enumi}{1}
\tightlist
\item
  Характеристическая функция случайной величины:
\end{enumerate}

\[
\varphi_X(t)=\mathbb E e^{itX}
=\int_{\mathbb R} e^{itx}\,P_X(dx).
\]

\begin{enumerate}
\def\labelenumi{\arabic{enumi}.}
\setcounter{enumi}{2}
\tightlist
\item
  Характеристическая функция случайного вектора:
\end{enumerate}

\[
\varphi_X(t)=\mathbb E e^{i\langle t,X\rangle}.
\]

\begin{enumerate}
\def\labelenumi{\arabic{enumi}.}
\setcounter{enumi}{3}
\tightlist
\item
  Конечномерная характеристическая функция процесса:
\end{enumerate}

\[
\varphi_{s_1,\ldots,s_n}(u_1,\ldots,u_n)
=
\mathbb E\exp\left(i\sum_{k=1}^n u_kX_{s_k}\right).
\]

\begin{enumerate}
\def\labelenumi{\arabic{enumi}.}
\setcounter{enumi}{4}
\tightlist
\item
  Базовые свойства:
\end{enumerate}

\[
\varphi(0)=1,\qquad |\varphi(t)|\le1,\qquad
\varphi(-t)=\overline{\varphi(t)}.
\]

\begin{enumerate}
\def\labelenumi{\arabic{enumi}.}
\setcounter{enumi}{5}
\tightlist
\item
  Независимые суммы:
\end{enumerate}

\[
X\perp Y
\quad\Rightarrow\quad
\varphi_{X+Y}(t)=\varphi_X(t)\varphi_Y(t).
\]

\begin{enumerate}
\def\labelenumi{\arabic{enumi}.}
\setcounter{enumi}{6}
\tightlist
\item
  Формула обращения:
\end{enumerate}

\[
F(b)-F(a)
=
\lim_{T\to\infty}
\frac1{2\pi}
\int_{-T}^{T}
\frac{e^{-ita}-e^{-itb}}{it}\varphi(t)\,dt.
\]

\subsection{8. Типичные дополнительные
вопросы}\label{ux442ux438ux43fux438ux447ux43dux44bux435-ux434ux43eux43fux43eux43bux43dux438ux442ux435ux43bux44cux43dux44bux435-ux432ux43eux43fux440ux43eux441ux44b}

\textbf{Вопрос. Почему ограниченная измеримая функция интегрируема
именно на конечной мере?}

Потому что

\[
\int |f|\,d\mu\le M\mu(\Omega).
\]

Если \(\mu(\Omega)<\infty\), правая часть конечна. Если мера бесконечна,
оценка не дает конечности; пример \(f\equiv1\) на \(\mathbb R\)
показывает, что утверждение вообще неверно.

\textbf{Вопрос. Почему характеристическая функция всегда существует?}

Потому что комплексная экспонента имеет модуль \(1\):

\[
|e^{i\langle t,X\rangle}|=1.
\]

Значит это ограниченная случайная величина, а на вероятностном
пространстве ограниченные случайные величины интегрируемы.

\textbf{Вопрос. Чем характеристическая функция отличается от моментной
производящей функции?}

Характеристическая функция:

\[
\varphi_X(t)=\mathbb E e^{itX}
\]

существует для всех \(t\in\mathbb R\). Моментная производящая функция:

\[
M_X(t)=\mathbb E e^{tX}
\]

может не существовать ни в какой окрестности нуля. Например, у
распределения Коши характеристическая функция есть, а моменты и
\(M_X(t)\) не существуют.

\textbf{Вопрос. Что значит, что характеристическая функция -
преобразование Фурье распределения?}

Это значит, что вместо плотности или меры \(P_X\) рассматривается
интеграл

\[
\int e^{itx}\,P_X(dx).
\]

Если есть плотность \(p\), то это обычное преобразование Фурье:

\[
\varphi(t)=\int e^{itx}p(x)\,dx.
\]

\textbf{Вопрос. Можно ли из характеристической функции восстановить
распределение?}

Да. По теореме единственности характеристическая функция однозначно
задает распределение. В одномерном случае есть формула обращения для
\(F(b)-F(a)\) в точках непрерывности \(F\).

\textbf{Вопрос. Почему для суммы независимых величин характеристические
функции перемножаются?}

Потому что

\[
e^{it(X+Y)}=e^{itX}e^{itY},
\]

а для независимых случайных величин ожидание произведения равно
произведению ожиданий.

\textbf{Вопрос. Верно ли произведение характеристических функций без
независимости?}

Нет. Формула

\[
\varphi_{X+Y}(t)=\varphi_X(t)\varphi_Y(t)
\]

требует независимости. Без независимости распределение суммы зависит от
совместного распределения, а не только от маргинальных законов.

\textbf{Вопрос. Как записать характеристическую функцию конечномерного
распределения процесса?}

Для индексов \(s_1,\ldots,s_n\):

\[
\varphi_{s_1,\ldots,s_n}(u_1,\ldots,u_n)
=
\mathbb E\exp\left(i\sum_{k=1}^n u_kX_{s_k}\right).
\]

Это характеристическая функция вектора \((X_{s_1},\ldots,X_{s_n})\).

\textbf{Вопрос. Что дает положительная определенность характеристической
функции?}

Она является необходимым свойством любой характеристической функции:

\[
\sum_{j,k}c_j\overline{c_k}\varphi(t_j-t_k)\ge0.
\]

Интуитивно это следует из неотрицательности квадрата модуля. В обратную
сторону теорема Бохнера говорит, что непрерывная в нуле положительно
определенная функция с \(\varphi(0)=1\) является характеристической
функцией некоторого распределения.

\textbf{Вопрос. Как связаны характеристические функции и моменты?}

Если существует момент нужного порядка, то производные в нуле выражают
моменты:

\[
\varphi_X^{(k)}(0)=i^k\mathbb E X^k.
\]

Но существование \(\varphi\) само по себе не означает существование
моментов. Подробно это отдельная тема Билета 13.

\textbf{Вопрос. Как проверить независимость координат через
характеристическую функцию?}

Нужно проверить факторизацию:

\[
\varphi_{(X_1,\ldots,X_n)}(t_1,\ldots,t_n)
=
\prod_{k=1}^n\varphi_{X_k}(t_k).
\]

Если она выполнена для всех аргументов, координаты независимы.

\subsection{9. Типичные
ошибки}\label{ux442ux438ux43fux438ux447ux43dux44bux435-ux43eux448ux438ux431ux43aux438}

\begin{enumerate}
\def\labelenumi{\arabic{enumi}.}
\item
  Говорить, что ограниченная функция всегда интегрируема, не упоминая
  конечность меры.
\item
  Забывать измеримость: ограниченность без измеримости не дает интеграла
  Лебега как объекта на данном измеримом пространстве.
\item
  Путать интегрируемость \(X\) и существование характеристической
  функции. Для \(\varphi_X\) нужна интегрируемость \(e^{itX}\), а она
  есть всегда, потому что модуль равен \(1\).
\item
  Путать характеристическую функцию с моментной производящей функцией:
\end{enumerate}

\[
\mathbb E e^{itX}\ne \mathbb E e^{tX}.
\]

\begin{enumerate}
\def\labelenumi{\arabic{enumi}.}
\setcounter{enumi}{4}
\item
  Считать, что если нет математического ожидания, то нет
  характеристической функции.
\item
  Применять формулу
\end{enumerate}

\[
\varphi_{X+Y}=\varphi_X\varphi_Y
\]

без независимости.

\begin{enumerate}
\def\labelenumi{\arabic{enumi}.}
\setcounter{enumi}{6}
\tightlist
\item
  Считать, что плотность существует у любого распределения. Формула
\end{enumerate}

\[
\varphi(t)=\int e^{itx}p(x)\,dx
\]

верна только при наличии плотности.

\begin{enumerate}
\def\labelenumi{\arabic{enumi}.}
\setcounter{enumi}{7}
\tightlist
\item
  Забывать в векторном случае скалярное произведение:
\end{enumerate}

\[
e^{i\langle t,X\rangle},
\]

а не произведение \(e^{itX}\) в одномерном смысле.

\begin{enumerate}
\def\labelenumi{\arabic{enumi}.}
\setcounter{enumi}{8}
\item
  Писать формулу обращения без условия о точках непрерывности функции
  распределения.
\item
  Смешивать конечномерные распределения процесса с одномерными. Для
  процесса одномерные законы \(X_t\) сами по себе обычно не задают
  совместную зависимость между разными моментами времени.
\end{enumerate}

\subsection{10. Связи с другими
билетами}\label{ux441ux432ux44fux437ux438-ux441-ux434ux440ux443ux433ux438ux43cux438-ux431ux438ux43bux435ux442ux430ux43cux438}

\textbf{Билет 05.} Характеристическая функция определена для случайного
вектора как измеримого отображения. Измеримость нужна, чтобы
\(e^{i\langle t,X\rangle}\) была случайной величиной.

\textbf{Билет 06.} Распределение случайного вектора \(P_X\) - мера, по
которой интегрируется \(e^{i\langle t,x\rangle}\). Конечномерные
распределения процесса в этом билете получают характеристические
функции.

\textbf{Билет 07.} Интегрируемость и теорема Лебега о доминированной
сходимости нужны для существования и непрерывности характеристической
функции.

\textbf{Билет 09.} Независимость дает факторизацию совместного
распределения и произведение характеристических функций.

\textbf{Билет 11.} Формула замены переменной объясняет равенство

\[
\mathbb E e^{i\langle t,X\rangle}
=
\int e^{i\langle t,x\rangle}\,P_X(dx).
\]

\textbf{Билет 13.} Производные характеристической функции в нуле связаны
с математическим ожиданием, вторыми моментами и матрицей корреляции.

\textbf{Билет 15.} Согласованные конечномерные распределения можно
описывать через согласованные конечномерные характеристические функции.

\textbf{Билет 17.} Узкая стационарность означает неизменность всех
конечномерных распределений при сдвиге, а значит неизменность их
характеристических функций.

\textbf{Билет 18.} Гауссовские векторы и последовательности удобно
задаются характеристической функцией

\[
\exp\left(i\langle t,m\rangle-\frac12t^T\Sigma t\right).
\]

\textbf{Билет 20.} Характеристическая функция - Фурье-образ
распределения; это связывает вероятностные меры, дельта-функции и
обобщенные функции.

\textbf{Билет 21.} Метод характеристических функций используется в
доказательстве центральной предельной теоремы и закона больших чисел.

\subsection{11. Минимум на
удовлетворительно}\label{ux43cux438ux43dux438ux43cux443ux43c-ux43dux430-ux443ux434ux43eux432ux43bux435ux442ux432ux43eux440ux438ux442ux435ux43bux44cux43dux43e}

Нужно сказать три вещи.

\begin{enumerate}
\def\labelenumi{\arabic{enumi}.}
\tightlist
\item
  Если \(|f|\le M\) и \(\mu(\Omega)<\infty\), то
\end{enumerate}

\[
\int |f|\,d\mu\le M\mu(\Omega)<\infty.
\]

\begin{enumerate}
\def\labelenumi{\arabic{enumi}.}
\setcounter{enumi}{1}
\tightlist
\item
  Характеристическая функция случайной величины:
\end{enumerate}

\[
\varphi_X(t)=\mathbb E e^{itX}.
\]

Для вектора:

\[
\varphi_X(t)=\mathbb E e^{i\langle t,X\rangle}.
\]

\begin{enumerate}
\def\labelenumi{\arabic{enumi}.}
\setcounter{enumi}{2}
\tightlist
\item
  Она всегда существует, потому что \(|e^{i\langle t,X\rangle}|=1\), и
  однозначно задает распределение.
\end{enumerate}

Минимальный пример: для \(X\sim \operatorname{Bernoulli}(p)\)

\[
\varphi_X(t)=q+pe^{it}.
\]

\subsection{12. Минимум на
хорошо}\label{ux43cux438ux43dux438ux43cux443ux43c-ux43dux430-ux445ux43eux440ux43eux448ux43e}

Кроме минимума нужно назвать свойства:

\[
\varphi(0)=1,\qquad |\varphi(t)|\le1,\qquad
\varphi(-t)=\overline{\varphi(t)}.
\]

Также нужно объяснить:

\begin{itemize}
\tightlist
\item
  почему \(\varphi\) равномерно непрерывна;
\item
  почему для независимых сумм функции перемножаются;
\item
  как записывается характеристическая функция конечномерного
  распределения процесса;
\item
  что без конечности меры ограниченность не гарантирует интегрируемость.
\end{itemize}

\subsection{13. Что добавить для
отлично}\label{ux447ux442ux43e-ux434ux43eux431ux430ux432ux438ux442ux44c-ux434ux43bux44f-ux43eux442ux43bux438ux447ux43dux43e}

Для отличного ответа стоит добавить:

\begin{enumerate}
\def\labelenumi{\arabic{enumi}.}
\item
  Доказательство интегрируемости ограниченной функции через оценку
  \(M\mu(\Omega)\).
\item
  Доказательство существования характеристической функции через
  ограниченность комплексной экспоненты.
\item
  Положительную определенность:
\end{enumerate}

\[
\sum_{j,k} c_j\overline{c_k}\varphi(t_j-t_k)\ge0.
\]

\begin{enumerate}
\def\labelenumi{\arabic{enumi}.}
\setcounter{enumi}{3}
\item
  Теорему единственности и идею формулы обращения.
\item
  Факторизацию совместной характеристической функции как критерий
  независимости координат.
\item
  Согласованность конечномерных характеристических функций:
\end{enumerate}

\[
\varphi_{s_1,\ldots,s_n,s_{n+1}}(u_1,\ldots,u_n,0)
=
\varphi_{s_1,\ldots,s_n}(u_1,\ldots,u_n).
\]

\begin{enumerate}
\def\labelenumi{\arabic{enumi}.}
\setcounter{enumi}{6}
\tightlist
\item
  Предупреждение: характеристическая функция существует всегда, но
  моменты и производные могут не существовать.
\end{enumerate}

\subsection{14. Устная версия ответа на 5
минут}\label{ux443ux441ux442ux43dux430ux44f-ux432ux435ux440ux441ux438ux44f-ux43eux442ux432ux435ux442ux430-ux43dux430-5-ux43cux438ux43dux443ux442}

Сначала я сформулирую базовый факт об интегрируемости. Пусть
\((\Omega,\mathcal F,\mu)\) - пространство конечной меры, \(f\) измерима
и ограничена: \(|f|\le M\). Тогда

\[
\int |f|\,d\mu\le\int M\,d\mu=M\mu(\Omega)<\infty.
\]

Значит \(f\) интегрируема. В вероятности \(P(\Omega)=1\), поэтому всякая
ограниченная случайная величина интегрируема.

Теперь характеристические функции. Для случайной величины

\[
\varphi_X(t)=\mathbb E e^{itX}
=\int_{\mathbb R} e^{itx}\,P_X(dx).
\]

Для случайного вектора \(X\in\mathbb R^n\):

\[
\varphi_X(t)=\mathbb E e^{i\langle t,X\rangle}
=\int_{\mathbb R^n}e^{i\langle t,x\rangle}\,P_X(dx).
\]

Она всегда существует, потому что \(|e^{i\langle t,X\rangle}|=1\). Это
как раз применение первой части билета.

Основные свойства:

\[
\varphi(0)=1,\qquad |\varphi(t)|\le1,\qquad
\varphi(-t)=\overline{\varphi(t)}.
\]

Также \(\varphi\) равномерно непрерывна: разность оценивается через

\[
\mathbb E|e^{i\langle h,X\rangle}-1|,
\]

и это стремится к нулю по теореме Лебега о доминированной сходимости.

Если \(X\) и \(Y\) независимы, то

\[
\varphi_{X+Y}(t)=\mathbb E e^{itX}e^{itY}
=\varphi_X(t)\varphi_Y(t).
\]

Это одно из главных применений характеристических функций.

Для процесса \((X_s)\) конечномерная характеристическая функция имеет
вид

\[
\varphi_{s_1,\ldots,s_n}(u_1,\ldots,u_n)
=
\mathbb E\exp\left(i\sum_{k=1}^n u_kX_{s_k}\right).
\]

Она является характеристической функцией вектора
\((X_{s_1},\ldots,X_{s_n})\) и поэтому однозначно задает соответствующее
конечномерное распределение.

В конце можно привести пример. Для Бернулли:

\[
\varphi_X(t)=q+pe^{it}.
\]

Для нормального распределения:

\[
\varphi_X(t)=\exp\left(imt-\frac{\sigma^2t^2}{2}\right).
\]

Главная ловушка: характеристическая функция существует всегда, но это не
значит, что существуют математическое ожидание, дисперсия или моментная
производящая функция.

\subsection{15. Если осталось 2 минуты перед
экзаменом}\label{ux435ux441ux43bux438-ux43eux441ux442ux430ux43bux43eux441ux44c-2-ux43cux438ux43dux443ux442ux44b-ux43fux435ux440ux435ux434-ux44dux43aux437ux430ux43cux435ux43dux43eux43c}

Запомнить нужно следующий каркас.

\begin{enumerate}
\def\labelenumi{\arabic{enumi}.}
\tightlist
\item
  Ограниченная измеримая функция на конечной мере интегрируема:
\end{enumerate}

\[
|f|\le M,\quad \mu(\Omega)<\infty
\Rightarrow
\int |f|\,d\mu\le M\mu(\Omega).
\]

\begin{enumerate}
\def\labelenumi{\arabic{enumi}.}
\setcounter{enumi}{1}
\tightlist
\item
  Характеристическая функция:
\end{enumerate}

\[
\varphi_X(t)=\mathbb E e^{itX},
\qquad
\varphi_X(t)=\mathbb E e^{i\langle t,X\rangle}
\]

для вектора.

\begin{enumerate}
\def\labelenumi{\arabic{enumi}.}
\setcounter{enumi}{2}
\item
  Она всегда существует, потому что модуль экспоненты равен \(1\).
\item
  Главные свойства:
\end{enumerate}

\[
\varphi(0)=1,\qquad |\varphi(t)|\le1,\qquad
\varphi(-t)=\overline{\varphi(t)}.
\]

\begin{enumerate}
\def\labelenumi{\arabic{enumi}.}
\setcounter{enumi}{4}
\tightlist
\item
  Независимые суммы:
\end{enumerate}

\[
\varphi_{X+Y}=\varphi_X\varphi_Y.
\]

\begin{enumerate}
\def\labelenumi{\arabic{enumi}.}
\setcounter{enumi}{5}
\tightlist
\item
  Конечномерное распределение процесса:
\end{enumerate}

\[
\varphi_{s_1,\ldots,s_n}(u)
=
\mathbb E\exp\left(i\sum u_kX_{s_k}\right).
\]

\begin{enumerate}
\def\labelenumi{\arabic{enumi}.}
\setcounter{enumi}{6}
\tightlist
\item
  Главная ошибка: не забыть условие конечности меры для интегрируемости
  и условие независимости для произведения характеристических функций.
\end{enumerate}

\newpage

\end{document}
